\honest{Similarity Clusters}{Eliezer Yudkowsky}{2008}

Plato's Academy defined man as a ``featherless biped.'' Diogenes the Cynic ``is said to have'' produced a plucked chicken and declared ``Here is Plato's man,'' and the Platonists added ``with broad nails.'' \nb{The story is told by Diogenes Laertius, in the third century AD. The division into feathered and featherless bipeds is in Plato's \textsc{Statesman}.}

No dictionary has ever listed all that humans have in common. Red blood, five fingers and 23 pairs of chromosomes are shared with other species; fission reactors single out only humans, but not all humans. The right gene sequences might pick out all humans and only humans, but would still be far from all they share. Still, ``Look for featherless bipeds'' picks out some humans rather than houses or sandwiches. From those examples we learn what else they share: language, tools, red blood, death by hemlock. So the category grows richer, and when Diogenes shows his plucked chicken, ``we are not fooled.''

If Aristotelian logic were a good model of the mind, the Platonists would have said of the chicken, ``Yes, that's a human; what's your point?'' \nb{Aristotle's own rules for definitions reject one that fits something other than the thing defined: a definition ``ought to be peculiar'' to it. By those rules the chicken refutes the definition, as the Platonists concluded.}

A definition only has to lead you to the ``similarity cluster,'' ``the group of things that have many characteristics in common.'' Then it has done its job, and I can go on to say new things, such as that humans are currently mortal. If the first featherless biped you meet is a plucked chicken, I can amend the map to ``broad nails'' and add: see Diogenes there? He is human, I am, you are; that chimpanzee is not, though close. \nb{Wittgenstein described categories of this kind in 1953 as ``family resemblances,'' ``a complicated network of similarities overlapping and criss-crossing.'' The post published the same day, ``Typicality and Asymmetrical Similarity,'' cites the psychological research on the same idea.} So a dictionary is not a book of Aristotelian class definitions but ``a book of hints for matching verbal labels to similarity clusters,'' or to properties useful for telling clusters apart.
